\section{Introduction}
\label{sec:introduction}

After a flood, earthquake, or wildfire, a responder may ask an unmanned aerial vehicle (UAV) to ``inspect the severely damaged building north of the road'' rather than provide a precise coordinate.
Executing such an instruction requires language grounding, navigation, and damage assessment.
It also requires the agent to recognize when its first observation is insufficient.
From a bird's-eye view, small structures, domain shift, image registration errors, and visually subtle damage can make a confident one-shot judgment unsafe.

Vision-and-language navigation (VLN) is commonly studied from an embodied first-person viewpoint~\cite{anderson2018vision}.
Disaster mapping, in contrast, often uses georeferenced overhead imagery and paired pre/post-event observations, as in xBD/xView2~\cite{gupta2019xbd}.
We connect these settings through \emph{bird's-eye disaster VLN}: an agent receives a free-form instruction, navigates over a metric, georeferenced overhead scene, identifies a disaster-relevant target, and returns both an arrival decision and a damage judgment.
The formulation intentionally separates \emph{getting to the requested structure} from \emph{having enough evidence to judge it}.

The paper's main hypothesis is narrow: calibrated predictive uncertainty can be converted into an active sensing decision.
Rather than always trusting the first view or always descending, \ProjectName estimates uncertainty from a dual-temporal damage classifier, evaluates a deterministic information-gain proxy, and reinspects under explicit per-location and per-episode budgets.
This follows the active-perception principle that sensing and motion should be selected jointly~\cite{bajcsy1988active}.

The remaining system components make this decision executable.
A hierarchical semantic planner decomposes the instruction into ordered landmarks; a metric semantic map summarizes visible, explored, and candidate regions; a compact memory records prior observations; and standard perception modules provide object and region proposals.
These components are necessary infrastructure, not the empirical headline.
In particular, our present navigation runs do not show a statistically significant strict-SR gain from hierarchical planning, semantic mapping, memory, or reinspection.

Our current results support two limited claims.
First, under a frozen event-disjoint split, post-hoc temperature scaling improves in-domain ECE without changing top-1 accuracy, but the same temperature can worsen held-out ECE.
Second, difference-aware temporal fusion does not recover a large holdout accuracy advantage once training and evaluation events are disjoint.
The end-to-end reinspection result remains underpowered even after the strict evidence-rich rerun: absolute success stays low, pairwise differences are non-significant, and genuine triggers are still too rare for a superiority claim.

Our contributions are:
\begin{itemize}
    \item a task definition and benchmark protocol for language-guided navigation and damage judgment in georeferenced bird's-eye disaster imagery;
    \item a calibrated, budgeted active-reinspection policy that links dual-temporal uncertainty to UAV sensing actions;
    \item an honest empirical account of calibration, temporal fusion, and low-SR navigation results under a leakage-safe evaluation protocol.
\end{itemize}
